\documentclass[11pt]{article}
\usepackage[margin=1in]{geometry}
\usepackage{amsmath,amssymb,amsthm}
\usepackage{graphicx}
\usepackage{hyperref}

\newtheorem{theorem}{Theorem}
\newtheorem{definition}{Definition}
\newtheorem{remark}{Remark}

\newcommand{\F}{\mathbb F}
\newcommand{\R}{\mathbb R}
\newcommand{\C}{\mathbb C}
\newcommand{\PP}{\mathbb P}
\newcommand{\Gammahat}{\widehat{\Gamma}}
\newcommand{\codim}{\operatorname{codim}}
\newcommand{\spanop}{\operatorname{span}}

\title{Forward Orthogonal Chains and Hamel Dimension}
\author{FA/Banach automatic math run: \texttt{fa\_banach\_001}}
\date{2026-07-18}

\begin{document}
\maketitle

\section*{Status}

\textbf{Partial result likely valid.}  The source paper asks which normed-space
properties can be decided from Birkhoff--James orthogonality alone, and then
shows, among other things, that its projective di-orthograph detects finite
dimension, smoothness, and strict convexity in the stated settings
\cite{AKKRZ2021}.  The result below strengthens the dimension-recognition
part for smooth spaces: the graph determines not only whether the dimension is
finite, but the exact Hamel dimension cardinal.

\section*{Source Question}

The extracted question is Question 2.1 in arXiv:2109.09414, page 4:

\begin{center}
\includegraphics[width=0.95\linewidth]{figures/open_problem_crop.png}
\end{center}

The source immediately notes that the question is too broad for a definite
global answer and restricts to the dimension, smoothness, and rotundness.  The
present packet stays in that dimension-recognition lane.

\section*{Definitions}

Let \(X\) be a normed space over \(\F\in\{\R,\C\}\).  Recall that \(x\) is
Birkhoff--James orthogonal to \(y\), written \(x\perp y\), if
\[
  \|x+\lambda y\|\geq \|x\| \qquad(\lambda\in\F).
\]
The projective di-orthograph \(\Gammahat(X)\) has as vertices the one
dimensional subspaces \([x]=\F x\), \(x\neq 0\), and a directed edge
\([x]\to [y]\) exactly when \(x\perp y\).  In a smooth normed space every
nonzero \(x\) has a unique supporting functional \(f_x\) normalized by
\(\|f_x\|=1\) and \(f_x(x)=\|x\|\).  The James criterion recalled in the
source gives
\[
  x\perp y \quad\Longleftrightarrow\quad f_x(y)=0
\]
in the smooth case.

\begin{definition}
A \emph{forward orthogonal chain} in \(\Gammahat(X)\) is a well-ordered family
\((v_\alpha)_{\alpha<\kappa}\) of distinct vertices such that
\[
  v_\alpha\to v_\beta \qquad(\alpha<\beta<\kappa).
\]
Let
\[
  \operatorname{fo}(X)
  = \sup\{\kappa:\Gammahat(X)\text{ contains a forward orthogonal chain of
  cardinal }\kappa\}.
\]
\end{definition}

\section*{Proof Intuition}

In a smooth normed space, each vertex \(v=[x]\) carries a single hyperplane
\(\ker f_x\) of outgoing neighbors.  A forward chain therefore behaves like a
triangular biorthogonal system: earlier supporting functionals vanish on later
vectors.  This triangularity forces every finite subset of the chain to be
linearly independent, giving the upper bound by the Hamel dimension.

For the reverse inequality, run the same idea by transfinite recursion.  Before
stage \(\alpha<\dim X\), fewer than \(\dim X\) supporting hyperplanes have been
chosen, so their intersection is still nonzero.  Choosing the next vector in
that intersection makes all earlier vertices point to the new one.

\begin{theorem}
Let \(X\) be a smooth normed space over \(\F\in\{\R,\C\}\).  Then
\[
  \operatorname{fo}(X)=\dim_{\F} X .
\]
Consequently, among smooth normed spaces, the projective Birkhoff--James
di-orthograph determines the full Hamel dimension cardinal.
\end{theorem}

\begin{proof}
The zero-dimensional case is empty, so assume \(X\neq\{0\}\).  Put
\(\kappa=\dim_{\F}X\).

First let \((v_\alpha)_{\alpha<\lambda}\) be a forward orthogonal chain.  Choose
nonzero representatives \(x_\alpha\in v_\alpha\), and let \(f_\alpha\) be the
unique supporting functional at \(x_\alpha\).  If \(\alpha<\beta\), then
\(v_\alpha\to v_\beta\), hence \(x_\alpha\perp x_\beta\), and smoothness gives
\[
  f_\alpha(x_\beta)=0.
\]
We claim that the family \((x_\alpha)_{\alpha<\lambda}\) is linearly
independent.  Indeed, take a finite relation
\[
  a_1x_{\alpha_1}+\cdots+a_nx_{\alpha_n}=0,
  \qquad \alpha_1<\cdots<\alpha_n .
\]
Applying \(f_{\alpha_1}\) gives
\[
  a_1f_{\alpha_1}(x_{\alpha_1})=0,
\]
because \(f_{\alpha_1}(x_{\alpha_j})=0\) for \(j>1\).  Since
\(f_{\alpha_1}(x_{\alpha_1})=\|x_{\alpha_1}\|\neq 0\), we get \(a_1=0\).
Repeating the argument with \(f_{\alpha_2},\ldots,f_{\alpha_n}\) yields
all \(a_j=0\).  Thus every forward orthogonal chain has cardinal at most
\(\dim X\), so \(\operatorname{fo}(X)\leq\kappa\).

For the reverse inequality, well-order a Hamel basis cardinal of \(X\) as
\(\kappa\).  We construct by transfinite recursion vectors
\((x_\alpha)_{\alpha<\kappa}\) and their supporting functionals
\((f_\alpha)_{\alpha<\kappa}\) so that
\[
  f_\beta(x_\alpha)=0 \qquad(\beta<\alpha).
\]
Suppose the construction has been completed for all \(\beta<\alpha\).  Since
\(|\alpha|<\kappa\), the intersection
\[
  M_\alpha=\bigcap_{\beta<\alpha}\ker f_\beta
\]
has codimension at most \(|\alpha|<\kappa\).  Therefore \(M_\alpha\neq\{0\}\).
Choose \(0\neq x_\alpha\in M_\alpha\), and let \(f_\alpha\) be the unique
supporting functional at \(x_\alpha\).  Then for every \(\beta<\alpha\),
\(f_\beta(x_\alpha)=0\), so \(x_\beta\perp x_\alpha\).  Thus the vertices
\(([x_\alpha])_{\alpha<\kappa}\) form a forward orthogonal chain of cardinal
\(\kappa\).  Hence \(\operatorname{fo}(X)\geq\kappa\), completing the proof.
\end{proof}

\section*{Limitations and Novelty Check}

The theorem is not a full answer to Question 2.1 of the source paper; that
question is intentionally broad.  It also does not remove smoothness.  Without
smoothness, an edge \(x\to y\) may be witnessed by one supporting functional at
\(x\) and another edge \(x\to z\) by a different supporting functional, so the
triangular linear-independence argument above can fail.

The local run indexes were searched for arXiv:2109.09414, Birkhoff--James
orthogonality, Hamel dimension, forward chains, transitive cliques, and
dimension-cardinal language.  No exact duplicate was found.  A web search on
2026-07-18 found the later open-access follow-up
\emph{Birkhoff--James classification of norm's properties} \cite{GKSSZ2024}.
That paper advances finite-dimensional property recognition, including smooth
points, maximal faces, supremum norms, and Radon-plane examples, but I did not
find the forward-chain/Hamel-cardinal statement above.

\section*{Human Review Recommendation}

Review as a compact partial theorem.  The key point is the use of smoothness in
both directions: it converts graph edges into equations against unique support
functionals.  A possible follow-up would be to determine whether a modified
graph invariant can recover Hamel cardinality without smoothness.

\begin{thebibliography}{9}

\bibitem{AKKRZ2021}
Lj. Arambasic, A. Guterman, B. Kuzma, R. Rajic, and S. Zhilina.
\emph{What does Birkhoff--James orthogonality know about the norm?}
arXiv:2109.09414, 2021.

\bibitem{GKSSZ2024}
A. Guterman, B. Kuzma, S. Singla, and S. Zhilina.
\emph{Birkhoff--James classification of norm's properties.}
Advances in Operator Theory 9, Article 43, 2024.
\url{https://doi.org/10.1007/s43036-024-00321-0}.

\end{thebibliography}

\end{document}
